\chapter{Fordítás és futtatás}

\section{A dokumentum fordítása (LuaLaTeX)}

A \code{Docs} mappában lévő \code{Makefile} a következőket csinálja:

\begin{enumerate}
  \item lefuttatja a \code{scripts/generate\_stats.py} scriptet, ami a nyers CSV-kből előállítja a \code{generated/*.csv} fájlokat,
  \item majd kétszer meghívja a \code{lualatex}-et, hogy a tartalomjegyzék és hivatkozások rendben legyenek.
\end{enumerate}

Használat:

\begin{verbatim}
cd Docs
make
\end{verbatim}

A kimenet a \code{build/Docs.pdf} fájlba kerül.

\section{Natív könyvtárak (C) fordítása}

A natív és OpenCL-es komponensek MinGW alatt készültek. Példa:

\begin{verbatim}
cd aeslib
mingw32-make clean
mingw32-make all

cd ../crypto_aes_opencl
mingw32-make clean
mingw32-make all
\end{verbatim}

\textbf{Megjegyzés:} Az x86 buildhez 32 bites toolchain és 32 bites OpenCL runtime szükséges.
Ha a gépen csak 64 bites OpenCL ICD/runtime van, akkor a 32 bites processz nem fogja tudni betölteni
a szükséges DLL-eket (ez a PC1-en megfigyelt jelenséggel összhangban van).

\section{WinForms alkalmazás}

A WinForms rész Visual Studio / \code{dotnet build} segítségével fordítható.
A benchmarkok futtatásához érdemes:

\begin{itemize}
  \item release buildet használni,
  \item stabil power profilt beállítani (laptopon különösen),
  \item OpenCL esetén először warm-up-ot futtatni (a GUI-ben opció).
\end{itemize}
